\section{Pembangkitan Kunci Putaran (Key Expansion)}
\label{sec:key-expansion-aes}

Setiap putaran AES memakai kunci putaran (\textit{round key}) yang berbeda. Semua
kunci putaran dibangkitkan dari kunci eksternal pengguna yang disebut
\textit{cipher key} atau \textit{key}. Proses pembangkitannya dilakukan oleh fungsi
\textit{KeyExpansion()}.

Satuan kerja proses ini adalah \textit{word}, yaitu 32 bit atau 4 byte. Tinjau
sebuah larik \textit{key} yang panjangnya 16 byte dengan $N_r = 10$ putaran.
Kunci-kunci putaran disimpan di dalam larik \textit{rk}. Elemen awal \textit{key}
langsung menjadi $rk[0]$, sedangkan kunci putaran lainnya disimpan di
$rk[1], rk[2], \dots, rk[10]$.

\subsection{Algoritma KeyExpansion}

Untuk AES-128, algoritma pembangkitan kunci putaran adalah sebagai berikut.
\begin{enumerate}
    \item Salin elemen-elemen \textit{key} ke dalam larik $w[0], w[1], w[2], w[3]$.
    Larik $w[0]$ berisi empat elemen pertama \textit{key}, $w[1]$ berisi empat
    elemen berikutnya, dan seterusnya.
    \item Mulai dari $i = 4$ sampai $43$, lakukan langkah berikut.
    \begin{enumerate}
        \item Simpan $w[i-1]$ ke dalam peubah \textit{temp}.
        \item Jika $i$ kelipatan 4, jalankan fungsi $g$ terhadap \textit{temp}.
        \begin{itemize}
            \item \textbf{RotWord}: geser \textit{temp} satu byte ke kiri secara
            sirkuler.
            \item \textbf{SubWord}: substitusikan setiap byte hasil pergeseran
            dengan S-box.
            \item \textbf{XOR Rcon}: XOR hasilnya dengan konstanta putaran
            $Rcon[i/4]$.
        \end{itemize}
        \item XOR-kan $w[i-4]$ dengan \textit{temp}, lalu simpan hasilnya di
        $w[i]$.
    \end{enumerate}
\end{enumerate}

Nilai $Rcon$ berbeda untuk setiap $j = i/4$. Konstanta itu didefinisikan sebagai
$Rcon[j] = (RC[j], 0, 0, 0)$, dengan $RC[1] = 1$ dan
$RC[j] = \code{02} \bullet RC[j-1]$. Simbol $\bullet$ menyatakan perkalian di
$\text{GF}(2^8)$. Nilai $RC[j]$ dalam heksadesimal tercantum pada
Tabel~\ref{tab:rcon-aes}.

\begin{table}[htbp]
\centering
\caption{Konstanta putaran $RC[1]$ sampai $RC[10]$}
\label{tab:rcon-aes}
\small
\begin{tabular}{@{}l *{10}{c}@{}}
\hline
$j$ & 1 & 2 & 3 & 4 & 5 & 6 & 7 & 8 & 9 & 10 \\
\hline
$RC[j]$ & 01 & 02 & 04 & 08 & 10 & 20 & 40 & 80 & 1B & 36 \\
\hline
\end{tabular}
\sumber{Munir, \textit{IF4020 Kriptografi}, ITB.}
\end{table}

% Diagram rantai pembangkitan satu word kunci putaran AES.
% Dirujuk dari section-03.tex dengan % Diagram rantai pembangkitan satu word kunci putaran AES.
% Dirujuk dari section-03.tex dengan % Diagram rantai pembangkitan satu word kunci putaran AES.
% Dirujuk dari section-03.tex dengan \input{figures/key-expansion-aes}.
% Digambar sendiri karena diagram pada slide sumber berupa tangkapan layar
% yang resolusinya rendah dan labelnya bertumpuk.
\begin{figure}[htbp]
    \centering
    \begin{tikzpicture}[
        kotak/.style={draw, rounded corners=1pt, minimum height=0.8cm,
                      minimum width=1.45cm, font=\small, inner sep=2pt},
        xor/.style={draw, circle, inner sep=0.6pt, minimum size=0.62cm,
                    font=\small},
        panah/.style={-{Latex[length=2mm]}, thick},
    ]
        \node[kotak] (prev) at (0,0) {$w[i-1]$};
        \node[kotak] (rot)  at (2.05,0) {RotWord};
        \node[kotak] (sub)  at (3.95,0) {SubWord};
        \node[xor]   (x1)   at (5.65,0) {$\oplus$};
        \node[xor]   (x2)   at (7.35,0) {$\oplus$};
        \node[kotak, minimum width=1.75cm] (wi) at (9.35,0) {$w[i]$};

        \node[kotak, minimum width=1.75cm] (wkn) at (7.35,-1.95) {$w[i-N_k]$};
        \node[font=\small] (rcon) at (5.65,1.35) {$\text{Rcon}[i/N_k]$};

        \draw[panah] (prev) -- (rot);
        \draw[panah] (rot)  -- (sub);
        \draw[panah] (sub)  -- (x1);
        \draw[panah] (x1)   -- (x2);
        \draw[panah] (x2)   -- (wi);
        \draw[panah] (rcon) -- (x1);
        \draw[panah] (wkn)  -- (x2);

        % Kurung putus-putus menandai ketiga operasi yang tergabung dalam g().
        \draw[dashed, rounded corners=3pt]
            (1.05,-0.72) rectangle (6.32,0.72);
        \node[font=\small, anchor=west] at (1.05,-1.12) {$g(w[i-1])$};
    \end{tikzpicture}
    \caption{Rantai pembangkitan satu word kunci putaran pada AES. Bila $i$ kelipatan $N_k$, word sebelumnya dikenai fungsi $g$ lebih dahulu; bila tidak, ia di-XOR-kan langsung. Khusus AES-256, pada $i \bmod 8 = 4$ dijalankan SubWord tanpa RotWord dan tanpa Rcon}
    \label{fig:key-expansion-aes}
    \sumber{Munir, \textit{IF4020 Kriptografi}, ITB.}
\end{figure}
.
% Digambar sendiri karena diagram pada slide sumber berupa tangkapan layar
% yang resolusinya rendah dan labelnya bertumpuk.
\begin{figure}[htbp]
    \centering
    \begin{tikzpicture}[
        kotak/.style={draw, rounded corners=1pt, minimum height=0.8cm,
                      minimum width=1.45cm, font=\small, inner sep=2pt},
        xor/.style={draw, circle, inner sep=0.6pt, minimum size=0.62cm,
                    font=\small},
        panah/.style={-{Latex[length=2mm]}, thick},
    ]
        \node[kotak] (prev) at (0,0) {$w[i-1]$};
        \node[kotak] (rot)  at (2.05,0) {RotWord};
        \node[kotak] (sub)  at (3.95,0) {SubWord};
        \node[xor]   (x1)   at (5.65,0) {$\oplus$};
        \node[xor]   (x2)   at (7.35,0) {$\oplus$};
        \node[kotak, minimum width=1.75cm] (wi) at (9.35,0) {$w[i]$};

        \node[kotak, minimum width=1.75cm] (wkn) at (7.35,-1.95) {$w[i-N_k]$};
        \node[font=\small] (rcon) at (5.65,1.35) {$\text{Rcon}[i/N_k]$};

        \draw[panah] (prev) -- (rot);
        \draw[panah] (rot)  -- (sub);
        \draw[panah] (sub)  -- (x1);
        \draw[panah] (x1)   -- (x2);
        \draw[panah] (x2)   -- (wi);
        \draw[panah] (rcon) -- (x1);
        \draw[panah] (wkn)  -- (x2);

        % Kurung putus-putus menandai ketiga operasi yang tergabung dalam g().
        \draw[dashed, rounded corners=3pt]
            (1.05,-0.72) rectangle (6.32,0.72);
        \node[font=\small, anchor=west] at (1.05,-1.12) {$g(w[i-1])$};
    \end{tikzpicture}
    \caption{Rantai pembangkitan satu word kunci putaran pada AES. Bila $i$ kelipatan $N_k$, word sebelumnya dikenai fungsi $g$ lebih dahulu; bila tidak, ia di-XOR-kan langsung. Khusus AES-256, pada $i \bmod 8 = 4$ dijalankan SubWord tanpa RotWord dan tanpa Rcon}
    \label{fig:key-expansion-aes}
    \sumber{Munir, \textit{IF4020 Kriptografi}, ITB.}
\end{figure}
.
% Digambar sendiri karena diagram pada slide sumber berupa tangkapan layar
% yang resolusinya rendah dan labelnya bertumpuk.
\begin{figure}[htbp]
    \centering
    \begin{tikzpicture}[
        kotak/.style={draw, rounded corners=1pt, minimum height=0.8cm,
                      minimum width=1.45cm, font=\small, inner sep=2pt},
        xor/.style={draw, circle, inner sep=0.6pt, minimum size=0.62cm,
                    font=\small},
        panah/.style={-{Latex[length=2mm]}, thick},
    ]
        \node[kotak] (prev) at (0,0) {$w[i-1]$};
        \node[kotak] (rot)  at (2.05,0) {RotWord};
        \node[kotak] (sub)  at (3.95,0) {SubWord};
        \node[xor]   (x1)   at (5.65,0) {$\oplus$};
        \node[xor]   (x2)   at (7.35,0) {$\oplus$};
        \node[kotak, minimum width=1.75cm] (wi) at (9.35,0) {$w[i]$};

        \node[kotak, minimum width=1.75cm] (wkn) at (7.35,-1.95) {$w[i-N_k]$};
        \node[font=\small] (rcon) at (5.65,1.35) {$\text{Rcon}[i/N_k]$};

        \draw[panah] (prev) -- (rot);
        \draw[panah] (rot)  -- (sub);
        \draw[panah] (sub)  -- (x1);
        \draw[panah] (x1)   -- (x2);
        \draw[panah] (x2)   -- (wi);
        \draw[panah] (rcon) -- (x1);
        \draw[panah] (wkn)  -- (x2);

        % Kurung putus-putus menandai ketiga operasi yang tergabung dalam g().
        \draw[dashed, rounded corners=3pt]
            (1.05,-0.72) rectangle (6.32,0.72);
        \node[font=\small, anchor=west] at (1.05,-1.12) {$g(w[i-1])$};
    \end{tikzpicture}
    \caption{Rantai pembangkitan satu word kunci putaran pada AES. Bila $i$ kelipatan $N_k$, word sebelumnya dikenai fungsi $g$ lebih dahulu; bila tidak, ia di-XOR-kan langsung. Khusus AES-256, pada $i \bmod 8 = 4$ dijalankan SubWord tanpa RotWord dan tanpa Rcon}
    \label{fig:key-expansion-aes}
    \sumber{Munir, \textit{IF4020 Kriptografi}, ITB.}
\end{figure}


\subsection{Contoh Pembangkitan Kunci Putaran}

Misalkan \textit{key} dalam heksadesimal adalah
\[
\textit{key} = (\code{54}, \code{77}, \code{6F}, \code{20}, \code{4F}, \code{6E},
\code{65}, \code{20}, \code{4E}, \code{69}, \code{6E}, \code{65}, \code{20},
\code{54}, \code{77}, \code{6F}),
\]
yaitu teks ASCII ``Two One Nine Two''. Karena panjangnya 16 byte, $N_k = 4$ dan
$N_r = 10$.

Empat word pertama diambil langsung dari kunci:
\[
w[0] = (\code{54},\code{77},\code{6F},\code{20}), \quad
w[1] = (\code{4F},\code{6E},\code{65},\code{20}),
\]
\[
w[2] = (\code{4E},\code{69},\code{6E},\code{65}), \quad
w[3] = (\code{20},\code{54},\code{77},\code{6F}).
\]

Sekarang bangkitkan $w[4]$. Karena $i = 4$ kelipatan 4, fungsi $g$ dijalankan
terhadap $\textit{temp} = w[3]$.

\textbf{Langkah 1 --- RotWord.} Geser $w[3]$ satu byte ke kiri secara sirkuler:
\[
(\code{20},\code{54},\code{77},\code{6F})
\longrightarrow
(\code{54},\code{77},\code{6F},\code{20}).
\]

\textbf{Langkah 2 --- SubWord.} Baca S-box pada Tabel~\ref{tab:sbox-aes} untuk
setiap byte: $\text{S-box}(\code{54}) = \code{20}$,
$\text{S-box}(\code{77}) = \code{F5}$,
$\text{S-box}(\code{6F}) = \code{A8}$, dan
$\text{S-box}(\code{20}) = \code{B7}$. Jadi
\[
(\code{54},\code{77},\code{6F},\code{20})
\longrightarrow
(\code{20},\code{F5},\code{A8},\code{B7}).
\]

\textbf{Langkah 3 --- XOR Rcon.} Untuk $i = 4$ berlaku $j = 1$ dan
$Rcon[1] = (\code{01},\code{00},\code{00},\code{00})$, sehingga
\[
g(w[3]) = (\code{20},\code{F5},\code{A8},\code{B7}) \oplus
(\code{01},\code{00},\code{00},\code{00})
= (\code{21},\code{F5},\code{A8},\code{B7}).
\]

\textbf{Langkah 4 --- XOR dengan $w[i-4]$.} Hasilnya adalah
\[
w[4] = w[0] \oplus g(w[3])
= (\code{54},\code{77},\code{6F},\code{20}) \oplus
(\code{21},\code{F5},\code{A8},\code{B7})
= (\code{75},\code{82},\code{C7},\code{97}).
\]

Untuk $i = 5, 6, 7$ tidak ada fungsi $g$ yang dijalankan karena $i$ bukan
kelipatan 4. Cukup XOR dengan word yang berjarak empat:
\[
w[5] = w[4] \oplus w[1] = (\code{3A},\code{EC},\code{A2},\code{B7}),
\]
\[
w[6] = w[5] \oplus w[2] = (\code{74},\code{85},\code{CC},\code{D2}),
\]
\[
w[7] = w[6] \oplus w[3] = (\code{54},\code{D1},\code{BB},\code{BD}).
\]
% Koreksi: slide menulis w[3] = (54,77,6E,20) dan SubWord (20,F5,9F,B7).
% Byte terakhir w[3] adalah 6F, bukan 6E. Dengan S-box(6F) = A8, nilai yang
% benar adalah w[4] = 7582C797, w[5] = 3AECA2B7, w[6] = 7485CCD2,
% dan w[7] = 54D1BBBD. Semua nilai diperiksa dengan Python.

Keempat word itu menyusun kunci putaran pertama. Dituliskan sebagai matriks
\texttt{state}, dengan urutan kolom seperti pada Persamaan~\eqref{eq:aes-state},
\[
K_0 =
\begin{bmatrix}
\code{54} & \code{4F} & \code{4E} & \code{20} \\
\code{77} & \code{6E} & \code{69} & \code{54} \\
\code{6F} & \code{65} & \code{6E} & \code{77} \\
\code{20} & \code{20} & \code{65} & \code{6F}
\end{bmatrix},
\qquad
K_1 =
\begin{bmatrix}
\code{75} & \code{3A} & \code{74} & \code{54} \\
\code{82} & \code{EC} & \code{85} & \code{D1} \\
\code{C7} & \code{A2} & \code{CC} & \code{BB} \\
\code{97} & \code{B7} & \code{D2} & \code{BD}
\end{bmatrix}.
\]

\subsection{Kekhususan AES-256}

AES-256 berbeda dari AES-128 dalam tiga hal. Jumlah putarannya 14 kali, bukan 10.
Panjang kuncinya 256 bit, bukan 128 bit. Aturan \textit{expansion key}-nya juga
berbeda.

Pada AES-128, urutan \textit{expansion key} adalah RotWord, SubWord, XOR Rcon,
lalu XOR dengan word sebelumnya. Pada AES-256, ada operasi SubWord tambahan pada
word tertentu.
\begin{itemize}
    \item Jika $i \bmod 8 = 0$, jalankan RotWord, lalu SubWord, lalu XOR Rcon.
    \item Jika $i \bmod 8 = 4$, jalankan SubWord saja.
\end{itemize}
